\begin{lemma}
Let \(\varepsilon,\eta,\varepsilon_1,\varepsilon_2,\beta\in\mathbb R\) and
\(n_0\in\mathbb N\).  Assume
\[
0<\varepsilon,\qquad 0<\eta,\qquad \eta<\varepsilon,\qquad
\beta=\frac12,
\]
and
\[
\noderef{ParameterHierarchy}(\eta,\varepsilon_1,\varepsilon_2,n_0).
\]
Then, with high probability in \(n\) under the two-bite law
\[
\noderef{TwoBiteEventProbability}
\left(n,\noderef{TwoBiteNaturalReducedVertexCount}(n),
\noderef{TwoBiteEdgeProbability}(\beta,n),-\right),
\]
the final graph satisfies
\[
\noderef{IndependenceNumberLess}\!\left(
\noderef{TwoBiteFinalGraph}(\omega)_0,
(1+\varepsilon)\sqrt{(n:\mathbb R)\log(n:\mathbb R)}\right).
\]
\end{lemma}
\begin{proof}
Fix the parameters in the statement.  The equality \(\beta=1/2\) means that
the probability law in the target statement is the same as the law with edge
parameter \(1/2\).  For every \(n\), set
\[
M(n)=\noderef{TwoBiteNaturalReducedVertexCount}(n),\qquad
p(n)=\noderef{TwoBiteEdgeProbability}\!\left(\frac12,n\right),
\]
and define the common probability functional
\[
\mathbb P_n(E)
=\noderef{TwoBiteEventProbability}(n,M(n),p(n),E).
\tag{P}
\]
All probabilities below are taken with respect to this same \(\mathbb P_n\).
Also write
\[
x_\eta(n)=\noderef{TwoBiteIndependenceScale}(1+\eta,n)
=(1+\eta)\sqrt{(n:\mathbb R)\log(n:\mathbb R)},
\]
\[
x_\varepsilon(n)=\noderef{TwoBiteIndependenceScale}(1+\varepsilon,n)
=(1+\varepsilon)\sqrt{(n:\mathbb R)\log(n:\mathbb R)},
\]
and
\[
K_\eta(n)=\noderef{TwoBiteNaturalIndependenceScale}(1+\eta,n).
\]
The subscript \(\eta\) is important: all fixed-set estimates below use the
smaller scale parameter \(\eta\), while the final independence-number target is
the larger \(\varepsilon\)-scale \(x_\varepsilon(n)\).

By \(\noderef{NaturalIndependenceScaleFitsTarget}\), applied to
\(0<\eta<\varepsilon\), there is \(N_{\mathrm{fit}}\) such that for all
\(n\ge N_{\mathrm{fit}}\),
\[
(K_\eta(n):\mathbb R)<x_\varepsilon(n). \tag{1}
\]
By \(\noderef{NaturalParameterApproximation}\), for all sufficiently large \(n\)
the ceiling approximation gives
\[
x_\eta(n)\le (K_\eta(n):\mathbb R)<x_\eta(n)+1.
\tag{2}
\]
Since \(x_\eta(n)\to\infty\) and \((x_\eta(n)+1)/n\to0\), increase the threshold
once more to obtain \(N_{\mathrm{size}}\) such that for all
\(n\ge N_{\mathrm{size}}\),
\[
0<K_\eta(n)<n. \tag{3}
\]
Indeed, for such \(n\) we have \(1\le x_\eta(n)\) and \(x_\eta(n)+1<n\), and
(2) then forces \(1\le K_\eta(n)\) and \(K_\eta(n)<n\).  Hence
\(\texttt{Nat.choose}(n,K_\eta(n))\) is positive and counts the
\(K_\eta(n)\)-element subsets of the \(n\)-element vertex set \(\mathrm{Fin}(n)\).

Now define the three regularity components, all at the \(\eta\)-scale:
\[
\mathcal D_n(\omega)=\noderef{FiberAndDegreeEvent}(\omega,\varepsilon_2),
\]
\[
\mathcal M_n(\omega)=
\noderef{TwoBiteMediumClosedPairsEvent}(k=K_\eta(n),\eta,\varepsilon_1,\omega),
\]
\[
\mathcal S_n(\omega)=
\noderef{TwoBiteSmallClosedPairsEvent}(k=K_\eta(n),\eta,\varepsilon_1,\omega).
\]
Let
\[
\mathcal R_n=\mathcal D_n\cap\mathcal M_n\cap\mathcal S_n.
\]
Unfolding \(\noderef{TwoBiteRegularityEvent}\), this is exactly the event
\[
\noderef{TwoBiteRegularityEvent}
(k=K_\eta(n),\eta,\varepsilon_1,\varepsilon_2).
\tag{4}
\]

We first prove that \(\mathcal R_n\) holds with high probability under the
common law \(\mathbb P_n\).  Let \(\rho>0\).  The hierarchy
\(\noderef{ParameterHierarchy}(\eta,\varepsilon_1,\varepsilon_2,n_0)\) supplies
the positivity hypotheses needed by the three high-probability children.  Apply
\(\noderef{FiberAndDegreeConcentration}\) with \(\varepsilon_2\), with
\(\beta=1/2\), and with the function \(m(n)=M(n)\).  Its probability functional
is exactly \(\mathbb P_n\) by (P), so
\[
\mathbb P_n(\mathcal D_n)\to1.
\]
Since complements under the finite two-bite probability law satisfy
\(\mathbb P_n(E^c)=1-\mathbb P_n(E)\), choose \(N_D\) such that for
\(n\ge N_D\),
\[
\mathbb P_n(\mathcal D_n^c)<\rho/6. \tag{5}
\]

Apply \(\noderef{MediumClosedPairsBound}\) with the scale parameter
\(\eta\), the same \(\varepsilon_1,\varepsilon_2,n_0\), the same hierarchy
\(\noderef{ParameterHierarchy}(\eta,\varepsilon_1,\varepsilon_2,n_0)\),
\(\beta=1/2\), \(m(n)=M(n)\), and
\[
k(n)=K_\eta(n)=\noderef{TwoBiteNaturalIndependenceScale}(1+\eta,n).
\]
Again the probability functional is \(\mathbb P_n\), and the event in the child
is exactly \(\mathcal M_n\).  Thus \(\mathbb P_n(\mathcal M_n)\to1\), so choose
\(N_M\) such that for \(n\ge N_M\),
\[
\mathbb P_n(\mathcal M_n^c)<\rho/6. \tag{6}
\]

Apply \(\noderef{SmallClosedPairsBound}\) with the same data
\(\eta,\varepsilon_1,\varepsilon_2,n_0\), \(\beta=1/2\), \(m(n)=M(n)\), and
\[
k(n)=K_\eta(n)=\noderef{TwoBiteNaturalIndependenceScale}(1+\eta,n).
\]
Its law is still \(\mathbb P_n\), and its event is exactly \(\mathcal S_n\).
Choose \(N_S\) such that for \(n\ge N_S\),
\[
\mathbb P_n(\mathcal S_n^c)<\rho/6. \tag{7}
\]
For \(n\ge\max\{N_D,N_M,N_S\}\), the complement of the intersection satisfies
\[
\mathcal R_n^c
=\mathcal D_n^c\cup\mathcal M_n^c\cup\mathcal S_n^c.
\]
The finite union bound for \(\mathbb P_n\), followed by (5), (6), and (7), gives
\[
\mathbb P_n(\mathcal R_n^c)
\le
\mathbb P_n(\mathcal D_n^c)+\mathbb P_n(\mathcal M_n^c)+
\mathbb P_n(\mathcal S_n^c)
<\frac{\rho}{2}. \tag{8}
\]

Set \(\delta=\rho/2\).  Apply
\(\noderef{FixedSetIndependenceProbability}\) with this \(\delta>0\) and with
the hierarchy
\[
\noderef{ParameterHierarchy}(\eta,\varepsilon_1,\varepsilon_2,n_0).
\]
This is the point where the fixed-set bound uses \(\eta\), not
\(\varepsilon\).  The child gives a threshold \(N_{\mathrm{fix}}\ge n_0\) such
that whenever \(n>N_{\mathrm{fix}}\), \(m=M(n)\), \(p=p(n)\), and
\(I\subseteq\mathrm{Fin}(n)\) satisfies
\[
|I|=K_\eta(n)=\noderef{TwoBiteNaturalIndependenceScale}(1+\eta,n),
\]
we have
\[
\mathbb P_n\!\left(
I\text{ is independent in }\noderef{TwoBiteFinalGraph}(\omega)_0
\text{ and }\mathcal R_n(\omega)\right)
\le
\delta\,\binom{n}{K_\eta(n)}^{-1}. \tag{9}
\]
Indeed, the regularity event in the child is
\[
\noderef{TwoBiteRegularityEvent}
(k=|I|,\eta,\varepsilon_1,\varepsilon_2,\omega),
\]
and since \(|I|=K_\eta(n)\), this is the same event \(\mathcal R_n\) by (4).

Let \(G_n(\omega)=\noderef{TwoBiteFinalGraph}(\omega)_0\), and let
\(\mathcal B_n\) be the bad event that the desired final conclusion fails:
\[
\mathcal B_n(\omega)\Longleftrightarrow
\neg\noderef{IndependenceNumberLess}(G_n(\omega),x_\varepsilon(n)).
\]
Unfolding \(\noderef{IndependenceNumberLess}\), if \(\mathcal B_n(\omega)\)
holds, then there is a finite set \(J\subseteq\mathrm{Fin}(n)\) such that
\[
G_n(\omega)\text{ is independent on }J
\qquad\text{and}\qquad
x_\varepsilon(n)\le (|J|:\mathbb R). \tag{10}
\]
For \(n\ge N_{\mathrm{fit}}\), (1) and (10) give
\[
(K_\eta(n):\mathbb R)<x_\varepsilon(n)\le(|J|:\mathbb R).
\]
Thus \(K_\eta(n)\le |J|\) as natural numbers: if \(|J|<K_\eta(n)\), the
order-preserving cast \(\mathbb N\to\mathbb R\) would give
\((|J|:\mathbb R)<(K_\eta(n):\mathbb R)\), contradicting the display.
For \(n\ge N_{\mathrm{size}}\), (3) also gives \(0<K_\eta(n)\).  Since
\(K_\eta(n)\le |J|\), choose a \(K_\eta(n)\)-element subset
\[
I\subseteq J.
\]
For example, enumerate the finite set \(J\) and take any first
\(K_\eta(n)\) elements.  The set \(I\) is independent in \(G_n(\omega)\):
if an edge of \(G_n(\omega)\) had both endpoints in \(I\), then the same edge
would have both endpoints in \(J\), contradicting the independence of \(J\).
Therefore, on \(\mathcal B_n\cap\mathcal R_n\), there exists a
\(K_\eta(n)\)-element subset \(I\) which is independent and for which
\(\mathcal R_n\) holds.  Equivalently,
\[
\mathcal B_n\cap\mathcal R_n
\subseteq
\bigcup_{\substack{I\subseteq \mathrm{Fin}(n)\\ |I|=K_\eta(n)}}
\left\{
I\text{ is independent in }G_n
\right\}\cap\mathcal R_n. \tag{11}
\]

For \(n\ge N_{\mathrm{size}}\), the indexing family in (11) has exactly
\(\texttt{Nat.choose}(n,K_\eta(n))\) members, because \(0\le K_\eta(n)\le n\)
and \(\mathrm{Fin}(n)\) has \(n\) vertices.  Applying the finite union bound to
(11) and then using the fixed-set estimate (9) for every indexed set \(I\), we
obtain
\[
\mathbb P_n(\mathcal B_n\cap\mathcal R_n)
\le
\sum_{\substack{I\subseteq \mathrm{Fin}(n)\\ |I|=K_\eta(n)}}
\delta\,\binom{n}{K_\eta(n)}^{-1}.
\]
The sum has \(\binom{n}{K_\eta(n)}\) identical upper bounds, so
\[
\mathbb P_n(\mathcal B_n\cap\mathcal R_n)
\le
\binom{n}{K_\eta(n)}
\delta\,\binom{n}{K_\eta(n)}^{-1}
=\delta. \tag{12}
\]
The last equality uses the positivity of
\(\texttt{Nat.choose}(n,K_\eta(n))\) from (3), so multiplication by its real
reciprocal cancels.

Finally,
\[
\mathcal B_n\subseteq \mathcal R_n^c\cup(\mathcal B_n\cap\mathcal R_n).
\]
For all \(n\) above the maximum of the thresholds already chosen, (8) and (12)
give
\[
\mathbb P_n(\mathcal B_n)
\le
\mathbb P_n(\mathcal R_n^c)+\mathbb P_n(\mathcal B_n\cap\mathcal R_n)
<\frac{\rho}{2}+\delta
=\rho.
\]
Because \(\rho>0\) was arbitrary, \(\mathbb P_n(\mathcal B_n)\to0\).  Using the
same complement identity under \(\mathbb P_n\), the probability of
\[
\noderef{IndependenceNumberLess}(G_n(\omega),x_\varepsilon(n))
\]
tends to \(1\).  Since \(x_\varepsilon(n)=(1+\varepsilon)
\sqrt{(n:\mathbb R)\log(n:\mathbb R)}\), and since the original probability
law with \(\beta\) is the same law by \(\beta=1/2\), this is precisely the
asserted \(\noderef{WithHighProbability}\) statement.
\end{proof}
